{\Large\bfseries Mockup\par}
\bigskip

\fullname

30 minutes, one handwritten or 4x6 index card (or A6) cheat sheet allowed, no printed material.

\bigskip
\begin{center}\rule{0.4\linewidth}{0.4pt}\end{center}
\bigskip

\begin{question}
Consider the following code, where \lstinline|N| is a positive integer.

\begin{lstlisting}
/* dq is defined as below */
for (int i = 0; i < N; i++)
    dq.addFirst(i);
while (!dq.isEmpty())
    System.out.println(dq.removeLast());
\end{lstlisting}

Which definition of \lstinline|dq| gives the lowest worst-case number of operations for this code (select one)? Then write that complexity as a function of $N$.

\begin{itemize}[leftmargin=2em,itemsep=1pt,topsep=2pt]
\item[\cb{0}] \lstinline|Deque<Integer> dq = new ArrayDeque<>();|
\item[\cb{0}] \lstinline|Deque<Integer> dq = new LinkedList<>();|
\item[\cb{1}] both have the same complexity
\end{itemize}

Complexity: $O($\underline{\hspace{1in}}$)$ \ifanswers ($N$)\fi
\end{question}

\begin{question}
The following lines are the body of \lstinline|removeLast()| for a circular-array deque (field \lstinline|data|, \lstinline|front|, \lstinline|size|; method \lstinline|index(i)| returns \lstinline|(front + i) % data.length|), but they are in the wrong order.
The line \lstinline|return e;| is already placed correctly at the end; write the numbers 1 to 4 in the boxes for the remaining lines, to give the order in which they must be executed.

\medskip
\begin{tabular}{@{}l@{\quad}l@{}}
\ansbox{4} & \lstinline|size--;| \\[3pt]
\ansbox{2} & \lstinline|E e = data[last];| \\[3pt]
\ansbox{3} & \lstinline|data[last] = null;| \\[3pt]
\ansbox{1} & \lstinline|int last = index(size - 1);| \\
\end{tabular}

\medskip
\lstinline|return e;|
\end{question}

\begin{question}
What does the following code fragment print?

\begin{lstlisting}
Deque<Integer> in = new ArrayDeque<>();     // used as a stack
Deque<Integer> out = new ArrayDeque<>();    // used as a stack

in.push(8);
in.push(3);
if (out.isEmpty()) while (!in.isEmpty()) out.push(in.pop());
in.push(6);
System.out.print(out.pop() + " ");

in.push(1);
if (out.isEmpty()) while (!in.isEmpty()) out.push(in.pop());
System.out.print(out.pop() + " ");

if (out.isEmpty()) while (!in.isEmpty()) out.push(in.pop());
System.out.print(out.pop() + " " + out.pop());
\end{lstlisting}

Answer: \underline{\hspace{2in}} \ifanswers (8 3 6 1: elements are moved to \lstinline|out| only when it is empty, so the oldest element is always on top of \lstinline|out|; two stacks emulate a queue)\fi
\end{question}

\begin{question}
\lstinline|dq| is a \lstinline|Deque<Integer>|. Suppose you only ever add elements with \textbf{one} of its methods, and only ever remove them with \textbf{one} other method (always the same pair). For each pair below, mark whether it makes the deque behave like a \lstinline|Stack|, a \lstinline|Queue|, or neither.

\begin{center}
\begin{tabular}{@{}p{2.6in}ccc@{}}
\toprule
Add with / remove with & Stack & Queue & Neither \\
\midrule
\lstinline|push| / \lstinline|pop| & \cb{1} & \cb{0} & \cb{0} \\
\lstinline|offer| / \lstinline|poll| & \cb{0} & \cb{1} & \cb{0} \\
\lstinline|push| / \lstinline|pollLast| & \cb{0} & \cb{1} & \cb{0} \\
\lstinline|offerLast| / \lstinline|pollLast| & \cb{1} & \cb{0} & \cb{0} \\
\bottomrule
\end{tabular}
\end{center}
\end{question}

\begin{question}
Consider the following declarations.

\begin{lstlisting}
Queue<String> q = new LinkedList<>();
Deque<String> d = new ArrayDeque<>();
LinkedList<String> l = new LinkedList<>();
\end{lstlisting}

Mark the statements that compile (ignore what happens when they run).
\begin{itemize}[leftmargin=2em,itemsep=1pt,topsep=2pt]
\item[\cb{0}] \lstinline|q.get(0);|
\item[\cb{1}] \lstinline|l.get(0);|
\item[\cb{1}] \lstinline|d.pop();|
\item[\cb{0}] \lstinline|q.pop();|
\end{itemize}
\end{question}

\begin{question}
The following method of a \lstinline|MyArrayDeque| adds an element at the front, but one line contains a bug.
\lstinline|data| is the array, \lstinline|front| is the index of the first element, \lstinline|size| is the number of elements, \lstinline|grow()| doubles the capacity.
Which change fixes the method (select one)?

Hint: try it by hand on a deque where \lstinline|front| is already 0.

\begin{lstlisting}[numbers=left,numberstyle=\tiny,xleftmargin=2em]
public void addFirst(E e) {
    if (size == data.length) grow();
    front = (front - 1) % data.length;
    data[front] = e;
    size++;
}
\end{lstlisting}

\begin{itemize}[leftmargin=2em,itemsep=1pt,topsep=2pt]
\item[\cb{0}] line 2 becomes \lstinline|if (size == data.length - 1) grow();|
\item[\cb{1}] line 3 becomes \lstinline[keepspaces=true]|front = (front - 1 + data.length) % data.length;|
\item[\cb{0}] line 3 becomes \lstinline|front = Math.abs(front - 1) % data.length;|
\item[\cb{0}] line 4 becomes \lstinline[keepspaces=true]|data[front % data.length] = e;|
\item[\cb{0}] line 5 becomes \lstinline[keepspaces=true]|size = size + 1 % data.length;|
\end{itemize}
\end{question}

\begin{question}
\lstinline|dq| is an \lstinline|ArrayDeque<Integer>| with $N$ elements.

\begin{lstlisting}
boolean isPalindrome(Deque<Integer> dq) {
    while (dq.size() > 1) {
        if (!dq.removeFirst().equals(dq.removeLast()))
            return false;
    }
    return true;
}
\end{lstlisting}

\lstinline|isPalindrome| checks whether the sequence of elements reads the same forwards and backwards (and empties \lstinline|dq| in the process). What is the worst-case complexity of \lstinline|isPalindrome|, as a function of $N$?

Answer: \underline{\hspace{1in}} \ifanswers ($O(N)$: each iteration removes 2 elements in $O(1)$, so there are $O(N)$ iterations total)\fi
\end{question}

\begin{question}
(2 points) A standard queue implementation of \lstinline|Integer|s, based on a circular array, starts empty with an initial capacity of 4 (its capacity doubles when it grows). The following operations are performed in order:

\begin{lstlisting}
q.enqueue(1);
q.enqueue(2);
q.enqueue(3);
q.dequeue();
q.dequeue();
q.enqueue(4);
q.enqueue(5);
q.enqueue(6);
q.enqueue(7);
q.dequeue();
\end{lstlisting}

Draw the final state of the underlying array (every cell, whether it holds a reference to an \lstinline|Integer| object or none at all), the \lstinline|Integer| objects it references, and the final values of \lstinline|front| and \lstinline|size|.

Hint: the array stores references to \lstinline|Integer| objects on the heap, not the integers directly.

Hint: the cells that do not reference any object contain \lstinline|null|.

\ifanswers
\begin{verbatim}
array (cells 0..7): [null, ref->4, ref->5, ref->6, ref->7, null, null, null]
front = 1, size = 4
(before the grow: [ref->5, ref->6, ref->3, ref->4], front = 2; the grow
copies 3, 4, 5, 6 into cells 0..3 with front = 0, enqueue(7) goes to
cell 4, and the last dequeue() removes 3)
\end{verbatim}
\else
\vspace{2.5in}
\fi
\end{question}

\begin{question}
(2 points) Design a class \lstinline|TaskQueue| supporting \lstinline|enqueue(String name, String urgency)| (where \lstinline|urgency| is one of \lstinline|"Urgent"|, \lstinline|"Regular"|, \lstinline|"Relaxed"|) and \lstinline|dequeue()| (removes and returns the name of the next task to process).

\lstinline|dequeue()| must always return an \lstinline|"Urgent"| task if any are waiting; otherwise a \lstinline|"Regular"| one if any are waiting; otherwise a \lstinline|"Relaxed"| one. Among tasks of the same urgency, they are processed in the order they were enqueued (FIFO).

Describe your design, so that both \lstinline|enqueue| and \lstinline|dequeue| run in $O(1)$ time, regardless of how many tasks are waiting.

\ifanswers
\begin{verbatim}
Keep three plain queues, one per urgency tier -- e.g. an array
Queue<String>[] tiers = new Queue[3], indexed 0 = Urgent, 1 = Regular,
2 = Relaxed (or three separate fields urgent/regular/relaxed; same idea).

enqueue(name, urgency): offer name onto the one queue matching urgency.
                        O(1): a single offer.

dequeue(): check tiers[0], then tiers[1], then tiers[2], in that fixed
           order, and poll from the first one that is not empty.
           O(1): always exactly 3 isEmpty() checks and 1 poll, regardless
           of how many tasks are waiting in each tier.

Each tier preserves FIFO order on its own, and checking the tiers in a
fixed urgency order (not depending on N) is what keeps dequeue O(1)
instead of O(number of tiers with content) -- here that's a constant (3)
either way, but the key idea is: don't scan or sort the tasks themselves.
\end{verbatim}
\else
\vspace{2.5in}
\fi
\end{question}
